\section{Related Work}
\label{sec:related}

We place Related Work after the experiments for a reason given in the introduction: which threads of prior work are load-bearing depends on the calibration findings. Three threads are relevant.

\paragraph{Classical causal discovery from observations and interventions.} The constraint-based line begins with the PC algorithm \citep{spirtes2000causation} and the equivalence-class theory of \citet{verma1990equivalence} and \citet{andersson1997characterization}, which together establish that observational distributions identify the CPDAG but not the DAG. Score-based discovery was formalized by \citet{chickering2002ges}. The active side --- how many interventions suffice to orient the remaining edges, and how to choose them --- was put on formal footing by \citet{eberhardt2005interventions}, with tractable algorithms developed by \citet{hauser2012gies} under the name GIES (Greedy Interventional Equivalence Search). The PC and PC-plus-greedy baselines in our benchmark are instantiations of these lines; they set the reference point against which LLM agents must compete.

\paragraph{LLMs for causal reasoning.} A first wave of empirical work asked whether LLMs can answer causal queries at all. \citet{kiciman2023llmcausal} report that GPT-class models perform surprisingly well on pairwise causal-direction classification, sometimes rivaling classical methods. \citet{zecevic2023causalparrots} argue that such results reflect pattern-matching from training corpora (``causal parrots'') rather than causal reasoning and offer a framework for diagnosing the difference. \citet{jin2023cladder} build CLadder, a structured question-answering benchmark spanning Pearl's three rungs of causation, and find that LLMs struggle noticeably on interventional and counterfactual queries relative to associational ones. Closest to our setup, \citet{long2023llmgraphs} probe whether LLMs can reconstruct entire causal graphs from textual descriptions; they find the answer is brittle and prompt-sensitive. Our benchmark extends this thread from static question-answering and text-only graph construction to an interactive environment where the LLM must query data, spend an intervention budget, and commit to a DAG whose errors are scored by a formal contract.

\paragraph{Benchmarks and environments for LLM agents.} Recent benchmarks have moved LLM evaluation from static QA toward interactive environments with long-horizon action. SWE-bench \citep{jimenez2024swebench} requires an agent to patch real GitHub issues; AgentBench \citep{liu2024agentbench} aggregates eight interactive task families; WebArena \citep{zhou2024webarena} exercises web navigation; OSWorld \citep{xie2024osworld} extends this to a full computer-use setting; GAIA \citep{mialon2023gaia} stresses the conceptually-easy-but-operationally-hard gap between humans and agents. ACDB shares this design ethos --- paired seeds, deterministic scoring, a small and public action schema --- but targets a narrower axis: scientific reasoning under a formal causal model, where ground truth is defined by construction and the scoring contract decomposes cleanly into identifiable sub-tasks. The intended complement to this literature is a benchmark that trades generality for a setting in which classical methods provide a strong reference and the failure modes of LLMs can be named precisely.
